\chapter{Sagas, peces y bodas lésbicas}

\epigraf{Les lesbianes? Sempre amb educació, però que no em toquin la Constitució.}{Felipe VI}

\lettrine[loversize=0.2, lines=2]{T}{odo} empezó un 23 de abril, Diada de Sant Jordi, en el corazón de Caldes de Montbui. Las calles olían a rosas, libros nuevos y bocadillos de fuet. María, con su mochila llena de chapas feministas y una camiseta que decía “Athena is my queen”, rebuscaba entre las mesas de libros de segunda mano. Su objetivo: una edición decente de La Maldición del Titán.\\

Justo cuando estiró la mano para cogerla, otra mano se cruzó con la suya.\\

—¡Uy! Perdón — dijeron al unísono.\\

Se miraron. Cristina, con su pelo revuelto por el viento y una chaqueta con parches de queercore y runas griegas pintadas a mano, sonrió nerviosa.\\

—¿También fan de Percy Jackson? —preguntó María, con un brillo curioso en los ojos.\\

—¿Fan? Me he tatuado “what belongs to the sea will return to the sea” en la costilla —respondió Cristina.\\

Ahí, entre libros polvorientos y gente comprando rosas, nació algo. Empezaron a hablar. Primero sobre Percy. Luego sobre Annabeth. Luego sobre política queer, tortillas sin huevo y si los centauros pagarían alquiler o no. Siguieron hablando por Instagram, luego por Telegram, luego por notas de voz que duraban media hora. Tres meses después, un mensaje de Cristina:\\

> “¿Y si hacemos una locura?”\\

“¿Como adoptar otro gato?”\\

“No… casarnos.”\\

\saltescena

La boda fue en un huerto urbano de Terrassa, oficiada por su excéntrico pero leal amigo Grinder, que insistió en llevar una túnica de druida y decir: “Con el poder del WiFi y el consentimiento… os declaro esposas”.\\

Las invitadas aplaudieron, lloraron y bailaron al ritmo de una playlist que iba de Bad Gyal a Florence + the Machine sin pedir perdón.\\

\saltescena

Después de la boda, se mudaron juntas a Cerdanyola del Vallès, a un piso con vistas parciales a la montaña y muy completas a los contenedores de reciclaje.\\

Allí conocieron a sus vecinas: Kev y Lawre. Dos lesbianas con flow, energía de tía enrollada y una colección de plantas con nombres de diosas.\\

—¡Bienvenidas al bloque! —dijo Kev, entregándoles un libro con cinta morada.\\

Era Stone Butch Blues.\\

—Este libro te cambia la vida —añadió, guiñando un ojo.\\

Lawre apareció detrás con una pecera pequeña.\\

— Y estos son Ácido e Hialurónico. Dicen que los peces traen suerte y estabilidad emocional. Aunque bueno, eso depende del pez.\\

Ácido parecía tener una actitud punk. Hialurónico solo flotaba con estilo.\\

—Y cuando queráis —añadió Kev—, podéis venir con nosotras al Piso Warro. Es donde se junta nuestra gente. Croquetas, bolleras, personas no binarias, gamers, poetas, un DJ furry y una señora que hace tarot con cartas de Pokémon.\\

—Literalmente el sitio más guay de todo el Vallès —dijo Lawre, ya con la pecera en sus manos como si fuera un cáliz.\\

María y Cristina se miraron, con esa mezcla de amor tranquilo y emoción adolescente.\\

Habían empezado con Percy Jackson, pasado por una boda druídica y acabado con peces comunistas y vecinas queer mágicas.\\

Y lo mejor de todo: aún quedaba historia por vivir.

\setlength{\parindent}{20pt}